% GENERATED FILE. Edit canonical lessons, then run npm run generate:books.
% canonical-source-hash: fnv1a64:b4b6c3a68fc3129d
% canonical-lessons: IT-C393-carcere, IT-C393-detenuto, IT-C393-punire, IT-C393-punizione, IT-C393-sanzione

\chapter{Pagare per gli errori}
\label{ch:it-pagare-per-gli-errori}
\hlchaptermodality{393}

\begin{chapteropening}
\textbf{By the end of this chapter:} \emph{I can talk about punishment: prison, a prisoner, punishing someone, a punishment and a fine or sanction.}

Complete the last lesson of chapter 393: I can talk about punishment: prison, a prisoner, punishing someone, a punishment and a fine or sanction.
\end{chapteropening}

\section[il carcere]{il carcere — the prison, the jail}

\label{lesson:IT-C393-carcere}



\begin{quote}
Before the new one: say the Italian for the conviction, then the Italian for the penalty.
\end{quote}

\subsection*{You'll want to know: il carcere}
\textbf{il carcere} — \textquotedblleft{}the prison, the jail\textquotedblright{}.

\begin{quote}
Se ti prendono, vai in carcere.
\end{quote}

\begin{quote}
\emph{If they catch you, you go to prison.}
\end{quote}

Stress on the first syllable: \emph{càrcere}. The plural is feminine: \emph{le carceri}.

\begin{cousinweb}
\emph{carcere} is inherited from Latin \emph{carcer, carceris} \textquotedblleft{}a prison, an enclosure\textquotedblright{} (EtymDB); it also named the starting gates of a racecourse.
\end{cousinweb}

\begin{grammarlens}[title={What you've built}]
One more word for telling, explaining and describing at B1.
\end{grammarlens}

\subsection*{Your turn}
\begin{itemize}
\raggedright
  \item \emph{Say it:} \emph{il carcere}
  \item \emph{Say it:} \emph{il carcere}, once more
  \item \emph{Say it:} a sentence of your own with \emph{il carcere}
  \item \emph{Recall:} say the Italian for the conviction, then the Italian for the penalty, then say \emph{il carcere} again
\end{itemize}

\begin{tcolorbox}[breakable,colback=teal!4,colframe=teal!35!black,title={Before you move on}]
What does il carcere mean? (\textquotedblleft{}The prison, the jail\textquotedblright{}.) Say it once more.
\end{tcolorbox}
\section[il detenuto]{il detenuto — the prisoner, the inmate}

\label{lesson:IT-C393-detenuto}



\begin{quote}
Before the new one: say the Italian for the penalty, then the Italian for the prison.
\end{quote}

\subsection*{You'll want to know: il detenuto}
\textbf{il detenuto} — \textquotedblleft{}the prisoner, the inmate\textquotedblright{}.

\begin{quote}
I detenuti possono studiare in carcere.
\end{quote}

\begin{quote}
\emph{The inmates can study in prison.}
\end{quote}

Past participle of \emph{detenere} \textquotedblleft{}to hold\textquotedblright{}. A woman is \emph{la detenuta}.

\begin{cousinweb}
\emph{detenuto} is the past participle of \emph{detenere} \textquotedblleft{}to hold, to detain\textquotedblright{}, from Latin \emph{dētinēre} \textquotedblleft{}to hold back\textquotedblright{} (Lewis \& Short), built on \emph{tenēre} \textquotedblleft{}to hold\textquotedblright{}.
\end{cousinweb}

\begin{grammarlens}[title={What you've built}]
One more word for telling, explaining and describing at B1.
\end{grammarlens}

\subsection*{Your turn}
\begin{itemize}
\raggedright
  \item \emph{Say it:} \emph{il detenuto}
  \item \emph{Say it:} \emph{il detenuto}, once more
  \item \emph{Say it:} a sentence of your own with \emph{il detenuto}
  \item \emph{Recall:} say the Italian for the penalty, then the Italian for the prison, then say \emph{il detenuto} again
\end{itemize}

\begin{tcolorbox}[breakable,colback=teal!4,colframe=teal!35!black,title={Before you move on}]
What does il detenuto mean? (\textquotedblleft{}The prisoner, the inmate\textquotedblright{}.) Say it once more.
\end{tcolorbox}
\section[punire]{punire — to punish}

\label{lesson:IT-C393-punire}



\begin{quote}
Before the new one: say the Italian for the prison, then the Italian for the prisoner.
\end{quote}

\subsection*{You'll want to know: punire}
\textbf{punire} — \textquotedblleft{}to punish\textquotedblright{}.

\begin{quote}
Se rompi la regola, ti puniscono.
\end{quote}

\begin{quote}
\emph{If you break the rule, they will punish you.}
\end{quote}

Conjugated with \emph{-isc-}: \emph{punisco, punisci}.

\begin{cousinweb}
\emph{punire} is inherited from Latin \emph{pūnīre}, older \emph{poenīre}, \textquotedblleft{}to punish\textquotedblright{} (Lewis \& Short), built on \emph{poena} \textquotedblleft{}penalty\textquotedblright{}.
\end{cousinweb}

\begin{grammarlens}[title={What you've built}]
One more word for telling, explaining and describing at B1.
\end{grammarlens}

\subsection*{Your turn}
\begin{itemize}
\raggedright
  \item \emph{Say it:} \emph{punire}
  \item \emph{Say it:} \emph{punire}, once more
  \item \emph{Say it:} a sentence of your own with \emph{punire}
  \item \emph{Recall:} say the Italian for the prison, then the Italian for the prisoner, then say \emph{punire} again
\end{itemize}

\begin{tcolorbox}[breakable,colback=teal!4,colframe=teal!35!black,title={Before you move on}]
What does punire mean? (\textquotedblleft{}To punish\textquotedblright{}.) Say it once more.
\end{tcolorbox}
\section[la punizione]{la punizione — the punishment; the free kick}

\label{lesson:IT-C393-punizione}



\begin{quote}
Before the new one: say the Italian for the prisoner, then the Italian for to punish.
\end{quote}

\subsection*{You'll want to know: la punizione}
\textbf{la punizione} — \textquotedblleft{}the punishment; the free kick\textquotedblright{}.

\begin{quote}
Per punizione non guardi la televisione per una settimana.
\end{quote}

\begin{quote}
\emph{As a punishment you are not watching television for a week.}
\end{quote}

In football, \emph{un calcio di punizione} is a free kick.

\begin{cousinweb}
\emph{punizione} is a learned borrowing from Latin \emph{pūnītiō} \textquotedblleft{}a punishing\textquotedblright{} (Lewis \& Short), from \emph{pūnīre} \textquotedblleft{}to punish\textquotedblright{}, on \emph{poena} \textquotedblleft{}penalty\textquotedblright{}.
\end{cousinweb}

\begin{grammarlens}[title={What you've built}]
One more word for telling, explaining and describing at B1.
\end{grammarlens}

\subsection*{Your turn}
\begin{itemize}
\raggedright
  \item \emph{Say it:} \emph{la punizione}
  \item \emph{Say it:} \emph{la punizione}, once more
  \item \emph{Say it:} a sentence of your own with \emph{la punizione}
  \item \emph{Recall:} say the Italian for the prisoner, then the Italian for to punish, then say \emph{la punizione} again
\end{itemize}

\begin{tcolorbox}[breakable,colback=teal!4,colframe=teal!35!black,title={Before you move on}]
What does la punizione mean? (\textquotedblleft{}The punishment; the free kick\textquotedblright{}.) Say it once more.
\end{tcolorbox}
\section[la sanzione]{la sanzione — the sanction, the penalty, the fine}

\label{lesson:IT-C393-sanzione}



\begin{quote}
Before the new one: say the Italian for to punish, then the Italian for the punishment.
\end{quote}

\subsection*{You'll want to know: la sanzione}
\textbf{la sanzione} — \textquotedblleft{}the sanction, the penalty, the fine\textquotedblright{}.

\begin{quote}
Se paghi in ritardo, c'è una sanzione.
\end{quote}

\begin{quote}
\emph{If you pay late, there is a penalty.}
\end{quote}

In offices and on tax forms it often means a fine.

\begin{cousinweb}
\emph{sanzione} is a learned borrowing from Latin \emph{sānctiō} \textquotedblleft{}a decree; the clause that sets a penalty\textquotedblright{} (Lewis \& Short), from \emph{sancīre} \textquotedblleft{}to make sacred, to ratify\textquotedblright{}.
\end{cousinweb}

\begin{grammarlens}[title={What you've built}]
One more word for telling, explaining and describing at B1.
\end{grammarlens}

\subsection*{Your turn}
\begin{itemize}
\raggedright
  \item \emph{Say it:} \emph{la sanzione}
  \item \emph{Say it:} \emph{la sanzione}, once more
  \item \emph{Say it:} a sentence of your own with \emph{la sanzione}
  \item \emph{Recall:} say the Italian for to punish, then the Italian for the punishment, then say \emph{la sanzione} again
\end{itemize}

\begin{tcolorbox}[breakable,colback=teal!4,colframe=teal!35!black,title={Before you move on}]
What does la sanzione mean? (\textquotedblleft{}The sanction, the penalty, the fine\textquotedblright{}.) Say it once more.
\end{tcolorbox}
